% s-193-p193-m002; original PDF 193, 194
\Aphorism{s-193-p193-m002}{280}
Die vielbescholtene Unbescheidenheit des Genies, das sich selbst die Unsterblichkeit decretirt und im Angedenken der Menschheit fortzuleben wünscht, verstanden, bewundert, geliebt — ist sie nicht eigentlich ein Wunder von Bescheidenheit, von \Sperr{mangelndem} Selbstbewusstsein, von liebenswürdiger Einordnung in den Verband des „allgemein Menschlichen“? Das hohe Lied der Menschenverachtung, der wirkliche Tyrannenwahnsinn des bevorrechteten Geistes klingt ganz anders; aber ihr seid verwöhnt durch den Demokratismus der Schaffenden, die euch ihr Bestes geben und gönnen. Richard Wagner legt einmal — ihr werdet sagen, in einem Augenblick revolutionärer Unzurechnungsfähigkeit — das leidenschaftliche Bekenntniss ab, er wolle den Nibelungenring sich und seinen Freunden ein einziges Mal vorführen lassen und dann die Partitur verbrennen! Es könnte einen grimmigen Unschöpfer-Gott geben, der sein eigenes Werk immer wieder durch Sintfluth vernichtete, — einen ungetreuen Verwalter, der alle seine Geistesschätze ungenützt vermodern liesse, einen Künstler, der seine letzte Künstler-Wollust in dem Gedanken genösse, „dass Geschlecht auf Geschlecht geboren werden und sterben würde und dass die Grössten dieser Geschlechter ihr Leben einsetzen würden, um das zu erringen, was Er hätte geben können, wenn er nur die Hand hätte öffnen wollen!“ Erst ein solcher Künstler wäre unbescheiden, und, nach Kant, sogar unmoralisch; denn „als vernünftiges Wesen will er nothwendig, dass alle Vermögen in ihm entwickelt werden, weil sie ihm doch zu allerlei möglichen Absichten dienlich und gegeben sind“.
\par

